\section*{الفصل \ref{ch:b1:optical-instruments} --- الأجهزة البصرية والعين}

\begin{solution}{exo:b1:optical-instruments:1}
اللانهاية: $f' = \qty{17}{mm}$ و $V = 1/0.017 = \qty{58.8}{\delta}$.
وعلى \qty{25}{cm}: $V = 58.8 + 1/0.25 = \qty{62.8}{\delta}$ و $f' = \qty{15.9}{mm}$.
والسعة \qty{4.0}{\delta}.
\end{solution}

\begin{solution}{exo:b1:optical-instruments:2}
\omterm{def:b1:optical-instruments:eye}{العين} القصيرة البصر: يجب أن تُصوَّر اللانهاية عند \omterm{def:b1:optical-instruments:punctum}{النقطة البعيدة}، ومنه $f' = \qty{-40}{cm}$ و
$V = \qty{-2.5}{\delta}$. \omterm{def:b1:optical-instruments:eye}{والعين} الطويلة البصر: يجب أن يُصوَّر جسم على \qty{25}{cm}
تصويرًا تخيليًّا على \qty{1.0}{m}: $V = 1/\overline{OA'} - 1/\overline{OA}
= -1 + 4 = \qty{+3.0}{\delta}$.
\end{solution}

\begin{solution}{exo:b1:optical-instruments:3}
$G = 25/4.0 = 6.3$؛ والطابع في المستوي البؤري، على \qty{4.0}{cm} من
\omterm{def:b1:mirrors-thin-lenses:lens}{العدسة}. وتفصيل قدره \qty{2}{mm}: $\theta' = 2/40 = \qty{0.05}{rad}$، في مقابل
$2/250 = \qty{0.008}{rad}$ على \qty{25}{cm}.
\end{solution}

\begin{solution}{exo:b1:optical-instruments:4}
$G = 900/25 = 36$ و $900/9 = 100$؛ وطولا الأنبوب $f_1' + f_2' =
\qty{925}{mm}$ و\qty{909}{mm}؛ والقمر: $36 \times 0.52^\circ = 19^\circ$
و $52^\circ$.
\end{solution}

\begin{solution}{exo:b1:optical-instruments:5}
$\gamma_1 = -160/4 = -40$؛ و$G_2 = 250/20 = 12.5$؛ و$G = 500$.
و$\overline{F_1A} = -f_1'^2/\Delta = -16/160 = \qty{-0.10}{mm}$: أي الجسم على
\qty{4.1}{mm} من \omterm{def:b1:optical-instruments:microscope}{الشيئية}. ويجب أن يبقى الجسم على تلك المسافة
من \omterm{def:b1:optical-instruments:microscope}{الشيئية}، فيتحرك الأنبوب كله بمقدار \qty{10}{\micro m} ---
وهو مقبض الضبط الدقيق.
\end{solution}

\begin{solution}{exo:b1:optical-instruments:6}
\omterm{def:b1:optical-instruments:eye}{بالعين} المجرّدة: $1.5/3\times10^{-4} = \qty{5}{km}$. وبالمقراب: $\theta_{\min}
= 1.22 \times 5.5\times10^{-7}/0.100 = \qty{6.7e-6}{rad}$؛ فيمكن فصل
المصباحين حتى $1.5/6.7\times10^{-6} \approx \qty{220}{km}$ (من حيث
المبدأ --- إذ يتدخل الهواء وتكوّر الأرض قبل ذلك). و$G_r = 3\times10^{-4}/
6.7\times10^{-6} = 45$.
\end{solution}

\begin{solution}{exo:b1:optical-instruments:7}
$D = 50/2 = \qty{25}{mm}$. والتعريض $\propto N^2$: أي $(8/2)^2 = 16$ مرة
أطول، فيكون $16/500 \approx 1/30$~\unit{s}. وهي تشتري عمق ميدان (وزيوغًا
أقلّ) بثمن ضبابية الحركة أو حامل ثلاثي.
\end{solution}

\begin{solution}{exo:b1:optical-instruments:8}
$H = 35^2/(11 \times 0.03) = \qty{3.7}{m}$؛ والوضوح من $H/2 =
\qty{1.9}{m}$ إلى اللانهاية. وبالضبط المسبق على $H$ لا تحتاج الآلة إلى ضبط:
فتُلتقط الصورة لحظة وقوع المشهد.
\end{solution}

\begin{solution}{exo:b1:optical-instruments:9}
المنظار $8\times30$: \omterm{prop:b1:optical-instruments:telescope}{البؤبؤ الخارج} $30/8 = \qty{3.8}{mm} < \qty{7}{mm}$، فيدخل الضوء
كله إلى \omterm{def:b1:optical-instruments:eye}{العين}. والمنظار $7\times50$: \qty{7.1}{mm}، وهو مطابق للبؤبؤ الليلي،
مع ضوء مجموع أكثر $(50/30)^2 = 2.8$ مرة: فهو منظار الغسق.
أما $20\times30$: فبؤبؤه الخارج \qty{1.5}{mm}، وصورته خافتة، وحقله ضيق،
ورعشة اليد مكبَّرة عشرين مرة.
\end{solution}

\begin{solution}{exo:b1:optical-instruments:10}
\omterm{def:b1:mirrors-thin-lenses:lens}{العدسة} اللاصقة: $f' = \qty{-25}{cm}$ و $V = \qty{-4.0}{\delta}$. والنظارة:
\omterm{def:b1:optical-instruments:punctum}{النقطة البعيدة} على $25 - 1.5 = \qty{23.5}{cm}$ من \omterm{def:b1:mirrors-thin-lenses:lens}{العدسة}، ومنه $f' =
\qty{-23.5}{cm}$ و $V = \qty{-4.3}{\delta}$ --- أي أقوى. فعلى \omterm{def:b1:mirrors-thin-lenses:lens}{العدسة} أن
تكوّن الصورة عند \omterm{def:b1:optical-instruments:punctum}{النقطة البعيدة} \omterm{def:b1:optical-instruments:eye}{للعين}، وبعدُها عن \omterm{def:b1:mirrors-thin-lenses:lens}{العدسة}
يتوقف على موضع \omterm{def:b1:mirrors-thin-lenses:lens}{العدسة}.
\end{solution}

\begin{solution}{exo:b1:optical-instruments:11}
اللابؤرية: $F_1' = F_2$، والفصل $f_1' + f_2' = \qty{20}{cm}$ (وهو أقصر
من \qty{40}{cm} لمقراب كبلري له $G$ نفسه). و$G = -30/(-10) =
+3$: أي منتصبة. \omterm{prop:b1:optical-instruments:telescope}{والبؤبؤ الخارج} $\overline{O_2P'} = f_2'(f_1' + f_2')/f_1' =
-10 \times 20/30 = \qty{-6.7}{cm}$: أي تخيلي، داخل الأنبوب، حيث لا تستطيع
\omterm{def:b1:optical-instruments:eye}{العين} أن تكون؛ \omterm{def:b1:optical-instruments:eye}{والعين} خلف \omterm{def:b1:optical-instruments:microscope}{العينية} لا تلتقط إلا الحزم التي
تصادف بلوغها --- ومن هنا ضيق الحقل، وهو الشكوى الكلاسيكية من مناظير
الأوبرا.
\end{solution}

\begin{solution}{exo:b1:optical-instruments:12}
$(200/7)^2 = 820$؛ و$\Delta m = 2.5\log_{10}820 = 7.3$: فالقدر
الحدّي $13.3$. والتدفق $\propto 1/r^2$، ومنه يُرى المصباح
على بعد $\sqrt{820} = 29$ مرة.
\end{solution}

\begin{solution}{pb:b1:optical-instruments:1}
\textbf{1.} $F_2 = F_1'$: فالفصل $f_1' + f_2' = \qty{1225}{mm}$
(مع \omterm{def:b1:optical-instruments:microscope}{العينية} \qty{25}{mm}) أو \qty{1206}{mm} (مع \qty{6}{mm}).

\textbf{2.} من أجل $L_1$: $\overline{O_1A_1} = f_1'$ (فالصورة عند $F_1'$). ومن أجل $L_2$:
$\overline{O_2A_1} = -f_2'$، ومنه $1/\overline{O_2A'} = 1/f_2' + 1/(-f_2')
= 0$: أي صورة في اللانهاية.

\textbf{3.} تتجمّع الحزمة في المستوي البؤري المشترك على ارتفاع $h =
f_1'\alpha$ (بالشعاع المارّ بالمركز $O_1$)؛ ومن هناك يحدّد الشعاع المارّ بالمركز $O_2$
اتجاه الخروج، أي $\alpha' = -h/f_2' = -(f_1'/f_2')\alpha$.

\textbf{4.} $G = -48$ و $-200$؛ والقمر: $25^\circ$ و $104^\circ$.

\textbf{5.} لا معنى للأعلى والأسفل في السماء؛ أما منظار الرصد الأرضي
فيضيف زوج موشورات معدِّلة (أو ناقلًا عدسيًّا).

\textbf{6.} في المستوي البؤري للعدسة $L_1$، وقطرها $f_1'\theta = 1200 \times
9.1\times10^{-3} = \qty{11}{mm}$.

\textbf{7.} الجسم $O_1$ عند $\overline{O_2O_1} = -(f_1' + f_2')$:
$1/\overline{O_2P'} = 1/f_2' - 1/(f_1' + f_2') = f_1'/[f_2'(f_1' + f_2')]$؛
والتكبير $\overline{O_2P'}/\overline{O_2O_1} = -f_2'/f_1' = 1/G$،
والقطر $D/\abs G$.

\textbf{8.} مع \qty{25}{mm}: $\overline{O_2P'} = 25 \times 1225/1200 =
\qty{25.5}{mm}$، والقطر $150/48 = \qty{3.1}{mm}$. ومع \qty{6}{mm}:
\qty{6.0}{mm} خلفها، والقطر \qty{0.75}{mm}.

\textbf{9.} كل شعاع دخل \omterm{def:b1:optical-instruments:microscope}{الشيئية} يمرّ عبر \omterm{prop:b1:optical-instruments:telescope}{البؤبؤ
الخارج}: \omterm{def:b1:optical-instruments:eye}{فالعين} الموضوعة هناك تتلقى الحقل كله. أما أبعد منه
فلا تلتقط إلا الحزم المركزية: فينكمش الحقل إلى ثقب مفتاح.

\textbf{10.} $D/\abs G \leq \qty{7}{mm}$: ومنه $\abs G \geq 150/7 \approx 21$.

\textbf{11.} $\theta_{\min} = 1.22 \times 5.5\times10^{-7}/0.150 =
\qty{4.5e-6}{rad} = 0.92''$.

\textbf{12.} $4.5\times10^{-6} \times 3.84\times10^5 = \qty{1.7}{km}$؛
وبالعين المجرّدة $3\times10^{-4} \times 3.84\times10^5 = \qty{115}{km}$.

\textbf{13.} $G_r = 3\times10^{-4}/4.5\times10^{-6} = 67$؛ \omterm{def:b1:optical-instruments:microscope}{والعينية} \qty{25}{mm}
(أي $48$) أقربها إليه، ودونه قليلًا؛ أما \qty{6}{mm} (أي $200$) فتتجاوزه
ثلاث مرات.

\textbf{14.} عند $G = 200$ تُرى لطخة الحيود تحت $200 \times 4.5
\times10^{-6} = \qty{9e-4}{rad} > \epsilon$: أي ضبابية ظاهرة. كما يتوزع ضوء
القمر على مساحة أكبر $(200/48)^2 = 17$ مرة على
\omterm{def:b1:optical-instruments:eye}{الشبكية}: أي أخفت.

\textbf{15.} $1'' = \qty{4.85e-6}{rad}$: ومنه $D = 1.22\lambda/\theta =
\qty{140}{mm}$. وفوق ذلك تحدّ الرؤية الجوية التمييز من على الأرض
(ما لم تُستعمل بصريات تكيّفية)؛ ويبقى $D$ الأكبر يجمع ضوءًا أكثر،
فيبلغ أجرامًا أخفت.

\textbf{16.} $(150/7)^2 = 460$.

\textbf{17.} $\Delta m = 2.5\log_{10}460 = 6.7$: أي القدر $12.7$.

\textbf{18.} $P = 2\times10^{-3} \times \pi(0.075)^2 = \qty{3.5e-5}{W}$،
أي $460$ ضعف ما تجمعه \omterm{def:b1:optical-instruments:eye}{العين} المجرّدة. والصورة \omterm{def:b1:optical-instruments:eye}{الشبكية} أكبر مساحةً $\abs G^2$
مرة: فنسبة السطوع السطحي $460/G^2 = 0.20$ (عند $G = 48$) و
$0.012$ (عند $G = 200$) --- وهي تساوي $(D/\abs G)^2/d_p^2$، أي مربع نسبة
\omterm{prop:b1:optical-instruments:telescope}{البؤبؤ الخارج} إلى بؤبؤ \omterm{def:b1:optical-instruments:eye}{العين}.

\textbf{19.} من أجل جسم ممتدّ ينمو الضوء المجموع مثل $D^2$
بينما تنمو مساحة الصورة مثل $G^2 \geq (D/d_p)^2$: فلا يستطيع السطوع السطحي
في أحسن حال إلا أن يساوي قيمته \omterm{def:b1:optical-instruments:eye}{بالعين} المجرّدة (حين يكون \omterm{prop:b1:optical-instruments:telescope}{البؤبؤ الخارج} $=$ بؤبؤ \omterm{def:b1:optical-instruments:eye}{العين}). أما النجم فغير
مميَّز: إذ يقع ضوؤه كله في لطخة حيود واحدة، فينمو سطوعه
مثل $(D/d_p)^2$.

\textbf{20.} النقطة عند الزاوية $\beta$ تُصوَّر على ارتفاع $f_1'\beta$ في
المستوي البؤري؛ ويمرّر الحاجز ما يحقق $\abs{f_1'\beta} \leq \phi/2$: فالحقل
الكامل $\phi/f_1'$ (للزوايا الصغيرة).

\textbf{21.} $20/1200 = \qty{0.0167}{rad} = 0.95^\circ$: فيتّسع القمر كله.
أما $5/1200 = 0.24^\circ$: فلا يتّسع.

\textbf{22.} الحقل الظاهري $\phi/f_2'$: $20/25 = \qty{0.8}{rad} = 46^\circ$؛
و$5/6 = 48^\circ$.

\textbf{23.} $0.95^\circ = 3420''$: ومنه $3420/15 = \qty{230}{s}$، أي نحو أربع
دقائق؛ و$0.24^\circ = 864''$: أي \qty{58}{s}.

\textbf{24.} \omterm{def:b1:optical-instruments:microscope}{العينية} \qty{25}{mm} (أي $G = 48$، قرب $G_r$، ساطعة، والقمر
كله) من أجل القرص؛ و\qty{6}{mm} (أي $G = 200$) لنظرة أقرب إلى الفوهات،
أخفت وأضبب. وعدد الجهاز الوحيد: فوهات حتى
\qty{1.7}{km}.

\textbf{25.} مع $f_1'$ و $D$ نفسيهما: $G$ نفسه، والحقول \omterm{prop:b1:optical-instruments:telescope}{والبؤبؤ الخارج}
والتمييز والضوء نفسها (ناقصًا أيّ حجب مركزي). والمكاسب: لا
زيغ لوني. والتغيّر العملي: أن البؤرة تقع أمام
المرآة، فيجب أن تطوي مرآة ثانوية صغيرة الحزمةَ إلى خارج الأنبوب نحو
\omterm{def:b1:optical-instruments:microscope}{العينية}.
\end{solution}
